\documentclass[a4paper,11pt]{article}
\usepackage[utf8]{inputenc}
\usepackage[T1]{fontenc}
\usepackage[english]{babel}
\usepackage{graphicx}
\usepackage{hyperref}

\usepackage[utf8]{inputenc}
\usepackage{newunicodechar}
\newunicodechar{├}{|}      % or \ensuremath{\vert}
\newunicodechar{─}{-}      % a hyphen
\newunicodechar{│}{|}      % vertical bar
\newunicodechar{└}{'}      % maybe a quote

\title{Hierarchical Data Format Version 5 (HDF5)}
\author{Yann Chardon-Grossard}
\date{June 2025}

\begin{document}
\maketitle

\section*{Key Concepts}

\begin{quotation}
An HDF5 file is a container for two kinds of objects: \emph{datasets}, which are collections of data in array form, and \emph{groups}, which are folder‐like containers holding datasets and other groups.  Groups behave like dictionaries, and datasets like NumPy arrays.
\end{quotation}
\noindent\href{https://docs.h5py.org/en/stable/quick.html}{Reference: h5py Quick Start}.

\section*{What is HDF5?}
HDF5 (Hierarchical Data Format version 5) is a binary format optimized for storing large volumes of numerical data, together with an API (via \texttt{h5py} or \texttt{PyTables} in Python) that provides:
\begin{itemize}
  \item \emph{Hierarchical} organization into groups and datasets, similar to a file system;
  \item the ability to add \emph{attributes} (metadata) to any object;
  \item high‐performance \emph{random access} thanks to chunking and compression.
\end{itemize}

\section*{1. Hierarchical Model}
\begin{itemize}
  \item \textbf{Groups} (\texttt{Group}): folder‐like containers.
  \item \textbf{Datasets} (\texttt{Dataset}): multidimensional arrays (\texttt{numpy.ndarray}).
  \item \textbf{Attributes} (\texttt{attrs}): key–value pairs attached to a group or dataset.
\end{itemize}

\newpage
\medskip
\noindent Example structure:
\begin{verbatim}
/                             % root
├─ measurements/              % group
│  ├─ time             (dataset)
│  └─ temperature      (dataset)
├─ parameters/                % group
│  ├─ voltage          (dataset)
│  └─ current          (dataset)
└─ metadata                   % group
   ├─ experiment_date (attribute)
   └─ operator        (attribute)
\end{verbatim}

\section*{2. Why HDF5 Excels at Organizing Data}

HDF5 delivers clear and efficient data management thanks to several strengths:

\begin{enumerate}
  \item \textbf{Natural Organization}\\
    You arrange your data in a simple tree structure, like folders and files on your computer.  
    For example, \texttt{/measurements/temperature} or \texttt{/parameters/pressure}.  
    This lets you locate and navigate directly to the information you need.

  \item \textbf{Built‐in Metadata}\\
    Each group or dataset can carry its own labels (attributes), such as measurement units, creation date, or a description.  
    All relevant information stays together in one place, without external sidecar files.

  \item \textbf{Partial, High‐Performance Access}\\
    When your arrays are very large, HDF5 can split them into “chunks” and compress them.  
    You can read or write a single chunk without loading the entire dataset into memory.  
    The result: fast access, even for multi‐gigabyte datasets.

  \item \textbf{Interoperability}\\
    HDF5 is a standard format recognized by many tools and languages (C, Fortran, MATLAB, R, Julia, etc.).  
    You can share your file and open it elsewhere without complex conversion.

  \item \textbf{Scalability}\\
    Whether your data are a few megabytes or multiple terabytes, HDF5 scales.  
    Performance remains high for both data access and updates, thanks to its internal indexing and caching mechanisms.
\end{enumerate}

\end{document}
